\documentclass[11pt]{article}
\usepackage[margin=25mm]{geometry}
\usepackage{amsmath,amssymb,bm,booktabs,longtable,hyperref,graphicx}
\emergencystretch=3em
\title{Supplement: operator derivation and reproducibility for LBFlow}
\author{Yu Zhou\\Independent Researcher\\\texttt{zhouyu@live.com}}
\date{Research snapshot: 8 October 2026}
\begin{document}\maketitle
\section{Source identity and scope}
The historical baseline is DRLab 0.2.2.1, full revision:\\
\texttt{1bef9bbe7047ca74280f5870f4787be4df12fe20}.
The maintained LBFlow implementation contains the repairs; the original
archive is preserved. The main executable accepts one simulation configuration per MPI
job. A reentrant multi-object library interface is not verified.

Quantitative results use equal reference densities, unit lattice speed and
the standard pressure weights ($\alpha=4/9$ in D2Q9, $1/3$ in D3Q19).
SRT and MRT, relaxation times 0.6--1.4 and two arithmetic/harmonic
kinematic-viscosity choices are exercised. The physical force is force
density. An acceleration would require multiplying by density, which is not
the meaning of the constant \texttt{GV} value used here.

\section{Equilibrium and source moments}
For the actual lattice weights,
\begin{align}
 \sum_iw_i&=1,&\sum_iw_i e_{ia}&=0,&\sum_iw_i e_{ia}e_{ib}&=\delta_{ab}/3,\\
 \sum_iw_i e_{ia}e_{ib}e_{ic}e_{id}
 &= (\delta_{ab}\delta_{cd}+\delta_{ac}\delta_{bd}+\delta_{ad}\delta_{bc})/9.
\end{align}
Thus the equilibrium has mass $\rho$, momentum $\rho\bm u$, and second moment
$pI+\rho\bm u\bm u$, where
$p=\rho_0c_{p,0}^2+\rho_1c_{p,1}^2$.
The recovered phase weights yield
$c_{p,k}^2=3(1-\alpha_k)/5$ in D2Q9 and $(1-\alpha_k)/2$ in D3Q19.
The reference-density relation between the two alphas balances pure-phase
pressure; it does not establish thermodynamic consistency.

The unprefactored Guo source $G_i$ has
\begin{equation}
 \sum_iG_i=0,\quad \sum_i\bm e_iG_i=\bm F,\quad
 \sum_i\bm e_i\bm e_iG_i=\bm u\bm F+\bm F\bm u.
\end{equation}
These identities are tested against the source emitted by the production
function, at nonzero velocity and force in all coordinate components.
Actual $M$ and $M^{-1}$ are emitted and checked independently.

The force-integration counterexample follows directly from the original
shifted equilibrium. Its momentum increment is $s_j\tau F$.
The old D3Q19 momentum rates were zero; old D2Q9 momentum rates were one.
SRT's $s_j=1/\tau$ gives $F$ instead. The corrected source and half-force
velocity give net $F$ in both collision operators.

\section{Perturbation and recoloring}
The actual perturbation arrays satisfy
\begin{equation}
 \sum_iB_i=1/3,\qquad \sum_i\bm e_iB_i=0,\qquad
 \sum_i\bm e_i\bm e_iB_i=I/3.
\end{equation}
For a unit normal $\bm n$ and
$\Delta_i=9\sigma|\bm g|[w_i(\bm e_i\cdot\bm n)^2-B_i]/(4\tau)$,
\begin{equation}
 \sum_i\Delta_i=0,\quad \sum_i\bm e_i\Delta_i=0,\quad
 \sum_i\bm e_i\bm e_i\Delta_i=\frac{\sigma|\bm g|}{2\tau}(\bm n\bm n-I).
\end{equation}
The trace is nonzero, so an MRT implementation must account for energy and
other even moments as well as shear. Direct projection of the actual arrays
onto $M$, using 32 deterministic random normals, identifies coupled rows
\{1,2,7,8\} in D2Q9 and \{1,2,9,10,11,12,13,14,15\} in D3Q19
(zero-based indices). All use local rate $1/\tau$. Odd perturbation moments
vanish. Applying an additional Guo source prefactor to this calibrated
perturbation would change its surface tension and is not done.

Opposite directions have opposite recoloring cosine and equal equilibrium
weights at rest. Consequently the recoloring correction sums to zero, and
$f_i^{0,*}+f_i^{1,*}=f_i^{post}$ pointwise. Component momentum can change by
equal and opposite segregation fluxes; total momentum is unchanged.

If the original fixed MRT rate $1/\tau_0$ is combined with a perturbation
scaled by local $1/\tau(\chi)$, its capillary stress has factor
$\tau_0/\tau(\chi)$. A thin-interface estimate with arithmetic interpolation is
\begin{equation}
 \sigma_{eff}\simeq\sigma\tau_0
 \int_{-1}^{1}\frac{d\chi}{2\tau(\chi)}
 =\sigma\frac{\tau_0}{\tau_0-\tau_1}\log\frac{\tau_0}{\tau_1}.
\end{equation}
This is a leading continuum estimate, not an exact finite-grid prediction.
At prescribed $\sigma=0.005$, the original solver gives 0.0031483 for
$(\tau_0,\tau_1)=(0.6,1.4)$ after 40000 steps and 0.0076439 for the opposite
pair after 12000 steps. Relative late pressure changes are
$4.8\times10^{-5}$ and $3.2\times10^{-9}$. The corrected arithmetic cases have
Laplace errors 0.145\% and 0.170\%; the harmonic cases have 0.841\% and
1.159\%. Local relaxation therefore removes the major capillary calibration
defect without proving exact finite-grid viscosity independence.

\section{Time and checkpoint contract}
The integer-time incoming momentum is $\bm j_n=\rho\bm u_n-\bm F/2$.
At physical rest, populations are initialized as
$f_i^{eq}(0)-G_i(0)/2$. The normal collision, perturbation and recoloring
kernel is applied once to create the outgoing state at time zero.
For uniform constant forcing,
\begin{equation}
 \bm j_n=(n-1/2)\bm F,\quad \bm j_n^*=(n+1/2)\bm F,\quad
 \bm u_n=n\bm F/\rho.
\end{equation}
An outgoing checkpoint labelled $n$ is followed by streaming and reporting
at $n+1$. Its physical velocity, if reconstructed directly before streaming,
would use $(\bm j_n^*-\bm F/2)/\rho$, not the incoming-state plus sign.
The sidecar identifies the outgoing stage, lattice, dimensions, phase count,
time and next action. Population-file length is checked. The complete
configuration and geometry are recorded separately in run receipts; the
sidecar does not validate every physical parameter.

Nonfinite outgoing populations abort before checkpoint writing.
Convergence histories reset on restart, so the same later state is
reproduced while an identical automatic stopping time is not promised.
The native-double restart representation is tested on the same Linux
platform; cross-endian portability is not claimed.

\section{References and measurement protocols}
For a force-density-driven channel of wall spacing $H$,
\begin{equation}
 u(y)=\frac{F_x}{2\mu}y(H-y),\quad \overline u=\frac{F_xH^2}{12\mu}.
\end{equation}
Fluid-node coordinates are $y=1/2,3/2,\ldots,H-1/2$ relative to the lower
halfway wall. Force scales as $H^{-3}$ and duration as $H^2$, with fixed
lattice viscosity. This is diffusive refinement of the same physical channel.

For two layers, each solution is $-F_xy^2/(2\mu_k)+a_ky+b_k$.
Four linear conditions enforce two no-slip walls and continuity of velocity
and shear stress. A second reference solves the measured diffuse-viscosity
problem independently:
\begin{equation}
 C=F_x\frac{\int_0^H y/\mu(y)\,dy}{\int_0^H 1/\mu(y)\,dy},\quad
 u(y)=\int_0^y\frac{C-F_xs}{\mu(s)}\,ds.
\end{equation}
Viscosity is interpolated between node centers and numerical quadrature is
performed on 64 subdivisions per lattice spacing. This reference uses the
measured phase profile and is not a manufactured analytic interface profile.

Laplace pressure is compared with $(d-1)\sigma/R$. Subgrid intersections
on all Cartesian edges are fit with orthogonal-distance circle/sphere
regression. Pressure is averaged where $|\chi|>0.99$ and the sampling point
is at least three nodes from the fitted interface. Fit residuals, sample
counts, spurious velocity, component masses and late pressure changes are
retained. The smallest radii are finite-interface tests.

The wetting angle is measured through phase 1 using
\begin{equation}
 \theta=\arccos\frac{y_w-y_c}{R_c},\qquad y_w=0.5.
\end{equation}
The convention is consistent with $\bm n_\chi\cdot\bm n_w=\cos\theta$,
where $\bm n_w$ points away from solid and $\bm n_\chi$ toward phase 0.
Solid-to-fluid contour crossings are excluded. The principal fit uses
$\chi=0$ and wall exclusion $0.75W_{0.9}$; the sensitivity envelope uses
$\chi=-0.1,0,0.1$ and exclusions $0.5,0.75,1,1.5$ interface widths.
Invalid fits are recorded rather than replaced with invented angles.

For a flat wall, weighted stencil reconstruction solves
$A\bm g=\bm b$ with $A=\sum_{fluid}w_i\bm e_i\bm e_i$.
In physical $(x,y)$ coordinates, $A=\operatorname{diag}(5/18,1/6)$, so
the clipped estimator $3\bm b$ has factors $5/6$ and $1/2$.
This linear-field result does not ensure a nonlinear diffuse profile,
halfway-wall placement or surface-energy consistency. The wall normal and
angle rotation are held common in the ablation.

\section{Experiment inventory and numerical data}
The main verification suite contains 16 force cases, 24 channel cases,
8 continuous/parallel/restart executions and 4 invalid-configuration cases.
Six global-field comparisons are bitwise identical. Twelve layered-flow
executions cover two viscosity permutations, three heights and two
interpolation laws. Thirty distinct wall configurations cover two gradients,
three angles and five radius/width combinations. Larger cases use complete
population checkpoints or fresh initialization and up to eight MPI ranks.
An interrupted initial batch is recorded; its six partial large-radius
macro-output sets are excluded and rerun from initial geometry.

\begin{longtable}{llrrrrr}
\caption{Principal contact-angle measurements, through phase 1.}\\
\toprule
Method & $\beta$ & $R_{ini}$ & Target & Measured & Late change & $u_{max}$ \\
\midrule\endfirsthead
\toprule Method & $\beta$ & $R_{ini}$ & Target & Measured & Late change & $u_{max}$ \\
\midrule\endhead
legacy & 0.35 & 24 & 60 & 65.733 & -0.00235 & 8.58e-05 \\
renormalized & 0.35 & 24 & 60 & 62.554 & -0.00465 & 9.21e-05 \\
legacy & 0.35 & 24 & 90 & 91.203 & 0.00326 & 2.09e-05 \\
renormalized & 0.35 & 24 & 90 & 91.248 & 0.00328 & 5.82e-05 \\
legacy & 0.35 & 24 & 120 & 117.123 & 0.00013 & 2.20e-04 \\
renormalized & 0.35 & 24 & 120 & 120.437 & 0.00019 & 1.01e-04 \\
legacy & 0.35 & 48 & 60 & 64.833 & -0.00002 & 1.15e-04 \\
renormalized & 0.35 & 48 & 60 & 61.644 & -0.00005 & 6.24e-05 \\
legacy & 0.35 & 48 & 90 & 90.590 & 0.00000 & 1.43e-05 \\
renormalized & 0.35 & 48 & 90 & 90.611 & 0.00000 & 5.27e-05 \\
legacy & 0.35 & 48 & 120 & 116.718 & 0.00000 & 1.83e-04 \\
renormalized & 0.35 & 48 & 120 & 120.014 & 0.00000 & 5.24e-05 \\
legacy & 0.70 & 12 & 60 & 66.398 & -0.00004 & 1.45e-04 \\
renormalized & 0.70 & 12 & 60 & 62.745 & -0.00013 & 1.91e-04 \\
legacy & 0.70 & 12 & 90 & 92.163 & 0.00044 & 5.10e-05 \\
renormalized & 0.70 & 12 & 90 & 92.309 & 0.00044 & 1.22e-04 \\
legacy & 0.70 & 12 & 120 & 118.453 & 0.00657 & 2.91e-04 \\
renormalized & 0.70 & 12 & 120 & 122.276 & 0.00897 & 1.58e-04 \\
legacy & 0.70 & 24 & 60 & 64.909 & -0.00227 & 1.83e-04 \\
renormalized & 0.70 & 24 & 60 & 61.295 & -0.00530 & 1.38e-04 \\
legacy & 0.70 & 24 & 90 & 91.041 & 0.00205 & 3.27e-05 \\
renormalized & 0.70 & 24 & 90 & 91.109 & 0.00207 & 1.09e-04 \\
legacy & 0.70 & 24 & 120 & 117.594 & 0.00006 & 2.57e-04 \\
renormalized & 0.70 & 24 & 120 & 121.349 & 0.00010 & 1.11e-04 \\
legacy & 0.70 & 36 & 60 & 64.518 & -0.00017 & 1.94e-04 \\
renormalized & 0.70 & 36 & 60 & 60.909 & -0.00055 & 1.30e-04 \\
legacy & 0.70 & 36 & 90 & 90.684 & 0.00283 & 2.78e-05 \\
renormalized & 0.70 & 36 & 90 & 90.730 & 0.00292 & 1.05e-04 \\
legacy & 0.70 & 36 & 120 & 117.184 & 0.00000 & 2.42e-04 \\
renormalized & 0.70 & 36 & 120 & 120.919 & 0.00000 & 9.46e-05 \\
\bottomrule
\end{longtable}


The companion \texttt{data/measurements.json} and \texttt{data/wetting.csv}
contain fitted radii, sensitivity bounds, late-time changes and spurious
velocity. Raw VTI, parameter and log files reside in the experiment artifact
directories selected when running the scripts. Measurement summaries are
included with the source repository; raw VTI fields are not bundled. Scripts regenerate
geometries and figures. No stochastic error bars or confidence intervals
are inferred from these deterministic calculations.

\section{Scientific limits and publication status}
Only equal-density, low-speed flow in the documented parameter range is
verified. Curved-wall and three-dimensional wetting accuracy, contact-line
dynamics, hysteresis, high density ratio, arbitrary pressure weights, and
phase-specific Darcy permeability are not established. The two viscosity
rules represent alternative finite-interface closures. The geometrical
projection and weighted linear reconstruction are established ideas;
the ablation does not demonstrate superiority over all published wetting
schemes or a new bulk model.

The sole author is Yu Zhou, Independent Researcher; the software license is MIT.
Public repository access, an immutable release citation, funding and competing
interest declarations remain to be finalized before submission.
The manuscript is a research draft, not a submitted article or a peer-review
decision. AI-assisted inspection and drafting are disclosed in the main text.
\end{document}
